\section{系统设计与方法}

\subsection{系统全貌与职责划分}

Cloudra 将有机反应机理探索组织为化学假设、几何操作、量子化学计算与证据更新相互衔接的研究过程。任务输入给出反应物、研究目标、反应条件及必要的方法约束。研究智能体据此提出候选基元步骤与竞争通道，调用工具构造和修改三维结构，并通过计算检验候选。交付证据有向无环图（Delivery DAG）持续记录物种、步骤及其证据状态，使后续探索能够从已有结构和计算事实继续推进。

系统的职责分为三个层次。模型承担机理假设、通道选择、计算方案与失败后的研究决策；化学信息学和几何工具承担原子身份关联、结构测量、局部坐标操作与路径准备；计算后端和证据检查程序承担实际求解、原生结果解析及对象相关的科学核验。Coordinator 组织研究任务，计算智能体执行相应工作，Supervisor 对照任务要求审阅交付图及提交证据。Supervisor 的审阅意见与频率、路径等计算事实分别保留，科学证据由对应的计算产物提供。

“智能的克制”体现在这些职责的边界上。模型以反应位点、运动片段和目标构型表达研究意图，程序维护该意图与具体原子、坐标和计算文件之间的对应关系。每轮探索既生成新的候选，也更新已有主张的证据状态。对于共享同一起点的竞争通道，候选构型和过渡态证据分别关联到各自步骤；这种组织方式使模型能够比较不同假设，同时保留各项结论的具体计算依据。

\subsection{化学信息学与几何算子：将研究意图转化为可计算结构}

Cloudra 以分子图承载化学身份，以关联的三维构型承载空间位置。模型选择反应位点、运动片段和目标取向，程序执行原子对应、结构测量、坐标变换及候选筛选。各操作共同形成“读取结构—提出局部改变—检查所得构型”的循环，使研究意图能够落实为可继续计算的结构。

\subsubsection{原子身份、跨端点对应与反应中心}

分子表示保留元素、同位素、键级、形式电荷及已声明立体信息。无三维坐标的组分通过 ETKDGv3 嵌入 \textsuperscript{6,7}；多组分体系分别构建后再组装，已有三维坐标予以保留。反应物和产物的对应关系首先由完整原子标号确定；未标号时，在元素和同位素守恒的重原子双射中搜索，以如下目标排列候选：

\begin{equation}\label{eq:methods-1}
\begin{aligned}
J(\pi)={}&2(N_{\mathrm f}+N_{\mathrm b})+N_{\mathrm o}\\
&+4\Delta_H+8\Delta_q+8\Delta_r+2N_{\mathrm s}.
\end{aligned}
\end{equation}

其中，$\pi$ 将产物原子对应至反应物原子；$N_{\mathrm f}$、$N_{\mathrm b}$ 和 $N_{\mathrm o}$ 分别为映射后新增键、删除键及键级改变的数量；$\Delta_H$、$\Delta_q$、$\Delta_r$ 分别是对应原子的氢数、形式电荷和自由基电子数之差的绝对值之和；$N_{\mathrm s}$ 为可比较的已声明立体环境的不一致数。该目标优先保留局部化学环境，并允许通过图编辑表达反应。两端均有四面体标签时，比较先校正邻接次序的奇偶性；若邻接集合已改变，则跳过这项手性比较。

当各元素、同位素分组的置换数乘积不超过 $10^5$ 时，比较全部元素守恒候选；空间较大时，以元素匹配、允许键级不同且限制完整环匹配的最大公共子图（MCS）生成种子 \textsuperscript{7}，再依氢数、邻接度和输入次序补齐未匹配原子。后一路径属于有限候选搜索。具有相同编辑事实和评分的候选合并，近优候选保留歧义说明。选定重原子对应后，将其绑定至具体构型，优先保留显式氢标号及局部同位素对应，枚举剩余氢的对应。完整全原子映射用于重排产物，使端点具有相同原子顺序；图差分给出成断键原子对，并单独记录键级变化。

\subsubsection{结构测量与局部刚体操作}

结构感知将坐标转换为键长、键角、二面角、环构象和近接触。对坐标 $\mathbf x_i$，原子间距为 $d_{ij}=\|\mathbf x_j-\mathbf x_i\|$；以 $j$ 为顶点的键角由两条键向量的归一化点积求得。距离与参考键长、范德华接触尺度的比较形成几何提示，分子图保留已有身份。纯坐标输入在入口感知连接关系；后续局部坐标变换沿用该图，并重新测量变换前后的异常。

非环键两侧的默认片段选择由图连通性完成：暂时移除指定键 $a-b$，取包含 $b$ 的连通分量 $\mathcal F_b$。若切键后两端仍连通，则无法按该键唯一选侧，环骨架的运动区域需由显式原子组给定。记 $\mathcal F$ 为显式指定的运动原子组，或切键得到的 $\mathcal F_b$。选定片段绕轴旋转时，对每个 $i\in\mathcal F$ 应用 Rodrigues 公式：

\begin{equation}\label{eq:methods-2}
\begin{aligned}
\mathbf x'_i&=\mathbf a+\mathbf R(\hat{\mathbf u},\theta)(\mathbf x_i-\mathbf a),\\
\mathbf R&=\cos\theta\,\mathbf I+(1-\cos\theta)\hat{\mathbf u}\hat{\mathbf u}^{\mathsf T}
+\sin\theta\,[\hat{\mathbf u}]_\times .
\end{aligned}
\end{equation}

这里 $\mathbf a$ 为轴上一点，$\hat{\mathbf u}$ 为轴的单位方向，$\theta$ 为目标转角，$\mathbf I$ 为单位矩阵，$[\hat{\mathbf u}]_\times\mathbf v=\hat{\mathbf u}\times\mathbf v$。操作保持所选片段内部距离，改变其相对于其余结构的取向。指定键伸缩则将 $\mathcal F_b$ 整体平移：

\begin{equation}\label{eq:methods-3}
\mathbf x'_i=\mathbf x_i+(L-d_{ab})\frac{\mathbf x_b-\mathbf x_a}{d_{ab}},\qquad i\in\mathcal F_b,
\end{equation}

其中 $L$ 是目标键长，其余原子坐标保持原值。手性反转通过围绕中心分别旋转两条可分离支链、交换其方向实现；基团替换则删除指定侧片段，将带连接点的新基团沿原键方向对接，并显式更新分子图。这些操作提供具有明确局部含义的候选构型，随后返回实际距离和接触变化。比较优化前后结构时，系统在相同元素序列下对距离感知的连接集合求差，列出新增或消失的接触及其片内、片间归属；这份记录用于诊断几何变化，其化学含义仍需结合研究假设判断。

\subsubsection{分支拥挤描述与反应物预组织}

对通过杂化、环结构及取代基图不等价检查的候选轴，系统用衰减的分支描述提供局部拥挤线索：

\begin{equation}\label{eq:methods-4}
B_s=\sum_{k\in\mathcal T_s}c_k r_k^{\mathrm{vdW}}\rho^{\ell_k},\qquad
\rho=0.5,\quad 0\leq\ell_k\leq5.
\end{equation}

$\mathcal T_s$ 是从取代基根原子出发、排除相邻轴端原子的遍历集合；$\ell_k$ 是首次访问深度，$r_k^{\mathrm{vdW}}$ 为原子范德华半径。仅当 $k$ 是该分支的根原子且其图度数至少为 3 时，取 $c_k=1.5$，其余原子取 $c_k=1$；候选轴的拥挤描述为两端四条分支的 $B_s$ 之和。深度衰减突出轴附近的贡献；当前实现以四分支总分不低于 9.0 作为候选轴筛选的经验条件。这是拓扑启发式描述，旋转能垒与选择性仍需势能面计算确定。

在 RDKit 预组织路径中，需要组装的多组分体系先固定首片段，对其余片段实施确定性的斐波那契球面方向采样及绕轴扭转。令采样数为 $K$，第 $s$ 个方向为

\begin{equation}\label{eq:methods-5}
\begin{aligned}
\mathbf u_s&=(a_s\cos s\gamma,\ y_s,\ a_s\sin s\gamma),\\
a_s&=\sqrt{1-y_s^2},\quad y_s=1-\frac{2s}{K-1},\\
\gamma&=\pi(3-\sqrt5).
\end{aligned}
\end{equation}

其中 $s=0,\ldots,K-1$；片段先将参考轴转向 $\mathbf u_s$，再绕此方向扭转 $2\pi s/K$。默认 $K=32$，每组放置的片内几何保持原值，片间碰撞通过整片平移排除。已声明反应原子对及其距离窗口时，放置按最大窗口偏差排序：

\begin{equation}\label{eq:methods-6}
S(\mathbf X)=\max_{(i,j)\in\mathcal W}
\{0,\,l_{ij}-d_{ij}(\mathbf X),\,d_{ij}(\mathbf X)-u_{ij}\}.
\end{equation}

$\mathbf X$ 为全部坐标，$\mathcal W$ 为受引导原子对集合，$[l_{ij},u_{ij}]$ 为各对的允许距离区间；所有距离均在区间内时 $S=0$。程序保留偏差最小的放置，加入 MMFF94 距离约束后弛豫。显式引导窗口保留至预组织结束，输出另行核对距离、立体信息、碰撞和片间接触连通性，再以重原子 RMSD 去重；候选排序使用去除约束惩罚项的力场单点能。单组分则把距离窗口写入距离几何边界矩阵，平滑边界后嵌入并受约束弛豫。所得结构属于预组织种子。

\subsubsection{端点对齐与图像依赖对势插值}

完成全原子对应后，两个端点分别去除坐标均值。令 $\mathbf P$、$\mathbf R$ 为产物与反应物的中心化坐标矩阵，每行对应一个原子。Kabsch 对齐求解

\begin{equation}\label{eq:methods-7}
\begin{aligned}
&\min_{\mathbf Q^{\mathsf T}\mathbf Q=\mathbf I,\,\det\mathbf Q=1}
\|\mathbf P\mathbf Q-\mathbf R\|_F^2,\\
&\mathbf P^{\mathsf T}\mathbf R=\mathbf U\Sigma\mathbf V^{\mathsf T},\\
&\mathbf Q=\mathbf U\operatorname{diag}[1,1,\det(\mathbf U\mathbf V^{\mathsf T})]\mathbf V^{\mathsf T}.
\end{aligned}
\end{equation}

$\|\cdot\|_F$ 为全部矩阵元素平方和的平方根，$\mathbf U\Sigma\mathbf V^{\mathsf T}$ 为奇异值分解；行列式修正将变换限制为正旋转，保持手性。\textsuperscript{1} 以对齐端点的笛卡尔线性插值初始化中间图像。对第 $k$ 个图像，$\lambda_k=k/(M-1)$，$M$ 为含两端的图像数，原子对目标距离及标准图像依赖对势（image-dependent pair potential, IDPP）目标为 \textsuperscript{2}

\begin{equation}\label{eq:methods-8}
\begin{aligned}
d^*_{ij,k}&=(1-\lambda_k)d^R_{ij}+\lambda_k d^P_{ij},\\
V_k&=\sum_{i<j}\frac{[d_{ij,k}-d^*_{ij,k}]^2}{d_{ij,k}^{4}}.
\end{aligned}
\end{equation}

$d^R_{ij}$、$d^P_{ij}$ 为端点距离，$d_{ij,k}$ 为当前图像距离。逆四次方权重随当前距离更新，使短距离的偏离受到更大惩罚。默认标准模式固定两端，以改进切线 NEB 耦合中间图像，通过 MDMin 优化几何目标。另有显式轻量近似：将分母替换为固定的 $\max(d^*_{ij,k},0.1\,\text{Å})^4$，用 Adam 独立更新各中间图像。两种模式记录各自的数值终止状态。候选帧按变化最大的监测原子对跨过端点距离中点的位置选取；生成路径及数值收敛均属于几何准备结果，过渡态身份由后续优化、频率和连通性证据检验。

\subsection{机理证据 DAG 与科学校验}

\subsubsection{科学对象、证据关联与研究契约}

Cloudra 将待研究机理表示为有向无环图 $G=(V,E)$。节点 $v\in V$ 表示具有明确化学身份的物种状态，边 $e=(v_i,v_j)\in E$ 表示从起始物种到目标物种的一个基元步骤。过渡态作为步骤的计算证据绑定到边上；同一物种的结构、频率和能量证据则关联到相应节点。图因而同时表达“提出了什么机理”与“哪些计算支持该机理”，而非将计算任务列表直接视为反应机理。

每个图元素携带研究契约规定的阶段集合 $K$。默认机理步骤要求鞍点优化、频率确证和短程 IRC 方向证明；初猜仅提供候选结构，不替代上述证据。热化学校正等阶段是否必需由具体任务决定。证据引用记录来源类型及标识；文件送审进一步生成带内容摘要的归档快照，并检查快照及 IRC 配套文件是否仍与记录一致。工具输出还可转换为类型化事实 $f=(c,o,r,d)$，其中 $c$ 为检查项，$o\in\{+,-,?\}$ 分别表示支持、反驳和不确定，$r$ 为来源引用，$d$ 为频率计数、能量等解析细节。事实记录保留解析版本与目标归属，使同一结论的支持和反证能够并存。归档快照固定该次送审所依据的证据版本。

系统分别记录阶段覆盖与科学支持。阶段绑定回答“哪份结果被提交给哪个阶段”；具体科学检查回答“该结果是否支持这一主张”。若任务显式声明必需检查集合 $R$ 与可选检查集合 $O$，令 $F_t$ 为归属于目标 $t$ 的事实，则其支持条件为

\begin{equation}\label{eq:methods-9}
\begin{aligned}
S_t={}&\mathbf 1[R\ne\emptyset]\\
&\land\bigwedge_{c\in R}\mathbf 1[\exists f\in F_t:f.c=c,\ f.o=+]\\
&\land\mathbf 1[\not\exists f\in F_t:f.c\in R\cup O,\ f.o=-].
\end{aligned}
\end{equation}

缺少正证或出现相关反证时，结论保持未解决；空检查集合不自动获得支持。频谱可解析也仅是观测事实，是否符合极小值或鞍点要求仍由对象相关的科学检查确定。

\subsubsection{驻点曲率与反应模式诊断}

驻点曲率的物理依据是质量加权 Hessian。设 $U(\mathbf x)$ 为势能，$\mathbf x$ 为笛卡尔坐标，$M$ 为每个原子质量沿三个坐标重复排列的对角矩阵，则

\begin{equation}\label{eq:methods-10}
\begin{aligned}
H_{\mathrm{mw}}&=M^{-1/2}\nabla_{\mathbf x}^{2}U\,M^{-1/2},\\
H_{\mathrm{mw}}\mathbf v_k&=\lambda_k\mathbf v_k.
\end{aligned}
\end{equation}

在排除整体平移和转动的振动子空间中，负本征值对应虚角频率；正常极小值不含负曲率方向，一阶鞍点恰有一个。实际校验读取量子化学后端输出的有符号波数，按 $n_{\mathrm{imag}}=\sum_k\mathbf 1[\tilde\nu_k<0]$ 计数，其中 $\tilde\nu_k$ 以负值表示虚频。物种极小值要求零虚频，过渡态要求恰有一个虚频，不另设忽略小负频的容差。缺失频率保留为未知。系统分别核对原生正常终止和优化收敛，并检查频率计算与优化末帧属于同一几何；频率单点本身不能证明结构已经优化收敛。\textsuperscript{5}

单一负曲率方向仍不足以确定反应身份。系统同时报告虚频模式中反应中心原子对的运动。在可解析的模式位移非零时，对原子位移 $\mathbf u_i$ 及监视原子对集合 $W$，相对位移诊断量为

\begin{equation}\label{eq:methods-11}
\eta=\frac{\max_{(i,j)\in W}\|\mathbf u_i-\mathbf u_j\|}
{\max_{i<j}\|\mathbf u_i-\mathbf u_j\|}.
\end{equation}

该量位于 0 与 1 之间，描述反应中心相对于全局运动的参与程度，用于确定后续需要检查的原子对及其运动。位移缺失或扭转运动难以用原子对距离解释时，系统保留相应诊断状态；方向验收依赖真实路径与端点证据。

\subsubsection{IRC 方向与端点身份}

IRC 在理论上是自过渡态出发、在质量加权坐标中的最速下降路径。\textsuperscript{3,4} 记 $\mathbf q=M^{1/2}\mathbf x$，$s\ge0$ 为分别从鞍点两侧向外增长的弧长，则两支路径满足

\begin{equation}\label{eq:methods-12}
\begin{aligned}
\frac{d\mathbf q_{\pm}}{ds}&=-\frac{\nabla_{\mathbf q}U}{\|\nabla_{\mathbf q}U\|},\\
\mathbf q_{\mathrm{init},\pm}&\approx\mathbf q_{\mathrm{TS}}\pm\epsilon\mathbf v_{-}.
\end{aligned}
\end{equation}

第一行为连续路径的定义；第二行为离开驻点的数值起始近似，其中 $\mathbf v_-$ 为唯一负曲率方向，$\epsilon>0$ 为初始微小位移。具体路径积分由量子化学后端完成。Cloudra 首先检查 IRC 起始几何与已验证鞍点一致，再提取两个方向各自真实的最外侧帧。对于两端均以完整极小值表示的步骤，每一侧原生优化日志的起始几何必须匹配相应外帧，优化必须收敛，其末帧必须与零虚频计算的几何一致，并对应声明的反应物或产物身份。两侧允许正向或交换配对。只要这些关联成立，短程 IRC 无需自行完整收敛到极小值；几何区间一致性仅用于提示检查轨迹。分离产物采用独立分支：核对真实产物外帧的完整化学图、各组分的零虚频极小值，并保留反应物方向的几何区间检查；此时不宣称已经验证一个结合产物络合物的极小值。

\subsubsection{证据汇总与定量主张}

证据验收依次解析原生事实、核对结构关联、执行阶段判据并汇总契约缺口。令 $K_e$ 为步骤 $e$ 所要求的阶段集合，$P_{e,k}$ 表示阶段 $k$ 已有与该步骤化学身份一致、且通过相应科学检查的证据，则该步骤的必需证据条件可写为

\begin{equation}\label{eq:methods-13}
A_e=\bigwedge_{k\in K_e}P_{e,k}.
\end{equation}

整项机理交付还需覆盖声明的目标物种与步骤。缺失材料、未解决的反证及待审主张继续关联到具体目标。Supervisor 对照完整任务审阅机理假设和证据的对应关系，提出补算或修订要求；已有频率、路径或身份缺口仍需相应计算事实补齐。工作流程允许带理由地跳过阶段或声明阶段不适用，这些处置与已有证据的科学支持分别记录。

对电子活化能垒，定量主张进一步关联鞍点和声明参照态的原生能量：

\begin{equation}\label{eq:methods-14}
\Delta E_e^{\ddagger}=E_{\mathrm{TS},e}-E_{\mathrm{ref},e}.
\end{equation}

其中两项能量采用相同单位及任务规定的理论级别，参照态由对应研究任务定义。节点和步骤记录各自的能量来源及结构身份，使数值差能够回溯到实际计算。需要焓或自由能时，热化学校正及其温度、标准态按任务要求另行声明，并与纯电子能差区分。

\subsection{引导式诊断与探索恢复}

机理探索中的失败并不具有相同的科学含义。Cloudra 将运行异常、数值不收敛与科学证据不足分别反馈：前者涉及输入、执行环境或产物读取；第二类依据原生日志中的优化或自洽场收敛信息识别；第三类包括缺失频率、一阶鞍点条件不满足，以及 IRC 端点不能对应声明物种等情况。运行完成与科学验证因而保持独立；一次初猜失败也不被解释为反应不可达。几何区间异常或振动模式对应关系不明确等提示单独报告，其建议性质不因被反馈给模型而改变。

诊断信息围绕可修正的对象组织，包括发生阶段、已观察到的事实、相关结构或日志，以及需要检查的输入项。反应中心编号无效时，反馈同时给出当前结构的原子表和编号范围；端点结构与声明物种不符时，反馈指向该端结构来源或物种声明，避免用修改映射掩盖物种错误。缺失频率时，候选结构保留未验证状态；IRC 端点优化不支持目标物种时，反馈要求复核相应端点来源与身份。这些检查定位了需要修正或补证的对象，失败的化学原因仍结合研究假设判断。

反馈通过两条互补路径进入后续推理：即时检查结果随工具调用返回，并与该次调用一起保留；较长计算结束后，系统向负责该步骤的智能体发送阶段、判定理由、证据位置和后续检查建议。已声明且成功完成的计算由程序衔接后续阶段，无需模型逐项轮询；需要修改研究输入或判断证据含义时，再交由模型处理。交付检查还将开放缺口定位到具体节点、边与阶段；督导意见可补充检查理由和后续建议，但不能替代所缺的原生计算证据。

这一闭环可用下式表示：

\begin{equation}\label{eq:methods-15}
d_t=\Phi(y_t,G_t,C),\qquad
a_{t+1}=\pi_{\mathrm{LLM}}(H_t,d_t).
\end{equation}

其中，$y_t$ 为本次工具或计算的观察结果，$G_t$ 为当前证据图，$C$ 为任务要求；$\Phi$ 将程序检查结果组织为诊断反馈 $d_t$，模型结合已有研究记录 $H_t$ 选择下一步行动 $a_{t+1}$。该式以 $\Phi$ 表示程序的事实整理，以 $\pi_{\mathrm{LLM}}$ 表示模型的研究决策，二者通过具体诊断相衔接。

恢复时，模型在原研究流程中声明要修改的输入与理由，程序保留历史尝试和原始产物，并重新安排受改动影响的后续步骤。模型负责选择新的构型、初猜或可区分竞争假设的检查；编号转换、证据关联和状态记录由程序承担。系统不因失败而静默更换理论方法或参数，也不将修复调用成功等同于机理已经成立。

\section*{参考文献}
\begingroup\fontsize{8.5}{10.5}\selectfont
\par\noindent\hangindent=1.5em\hangafter=1 [1] Kabsch, W. A discussion of the solution for the best rotation to relate two sets of vectors. Acta Crystallographica Section A 1978, 34, 827–828. \url{https://doi.org/10.1107/S0567739478001680}\par\smallskip
\par\noindent\hangindent=1.5em\hangafter=1 [2] Smidstrup, S.; Pedersen, A.; Stokbro, K.; Jónsson, H. Improved initial guess for minimum energy path calculations. The Journal of Chemical Physics 2014, 140, 214106. \url{https://doi.org/10.1063/1.4878664}\par\smallskip
\par\noindent\hangindent=1.5em\hangafter=1 [3] Fukui, K. Formulation of the reaction coordinate. The Journal of Physical Chemistry 1970, 74, 4161–4163. \url{https://doi.org/10.1021/j100717a029}\par\smallskip
\par\noindent\hangindent=1.5em\hangafter=1 [4] Maeda, S.; Harabuchi, Y.; Ono, Y.; Taketsugu, T.; Morokuma, K. Intrinsic reaction coordinate: Calculation, bifurcation, and automated search. International Journal of Quantum Chemistry 2015, 115, 258–269. \url{https://doi.org/10.1002/qua.24757}\par\smallskip
\par\noindent\hangindent=1.5em\hangafter=1 [5] FACCTs. ORCA 6.1 Manual: Vibrational frequencies and Intrinsic reaction coordinate. \url{https://www.faccts.de/docs/orca/6.1/manual/contents/structurereactivity/frequencies.html}; \url{https://www.faccts.de/docs/orca/6.1/manual/contents/structurereactivity/irc.html}\par\smallskip
\par\noindent\hangindent=1.5em\hangafter=1 [6] Riniker, S.; Landrum, G. A. Better Informed Distance Geometry: Using What We Know To Improve Conformation Generation. Journal of Chemical Information and Modeling 2015, 55, 2562–2574. \url{https://doi.org/10.1021/acs.jcim.5b00654}\par\smallskip
\par\noindent\hangindent=1.5em\hangafter=1 [7] RDKit. Getting Started in Python: Conformer Generation and Maximum Common Substructure. \url{https://www.rdkit.org/docs/GettingStartedInPython.html}\par\smallskip
\endgroup
